\begin{proposition}[Proved: reduced bases with regular arithmetic fibres can fail]
\label{C:reduced_regular_fibres}
Let
\[
 k=\bigl(\mathbb Z_{(2)}[x]/(2x)\bigr)_{(2,x)}.
\]
This is a reduced noetherian local ring of characteristic zero.
Every integer localization has $K_0=\mathbb Z$. Its two nonempty
arithmetic fibres are
\[
 k\otimes\mathbb Q=\mathbb Q,
 \qquad k\otimes\mathbb F_2=\mathbb F_2[x]_{(x)},
\]
both regular. Nevertheless $\vartheta_k$ is not surjective in
degree two. More precisely the kernel of
$R_k(C_2)\to R_{\mathbb Q}(C_2)$ contains a direct sum of
countably infinitely many copies of $\mathbb Z/2$, all of which
vanish on restriction to every residue-field point of
$\operatorname{Spec}k$.
\end{proposition}
\begin{proof}
The ideal $(2x)$ in $\mathbb Z_{(2)}[x]$ is $(2)\cap(x)$:
if $xf$ is divisible by $2$, reduction modulo $2$ shows that $f$
is divisible by $2$. Both ideals are prime, so the quotient and
its localization are reduced. Locality, noetherianity, the displayed
fibres, and $K_0(k)=\mathbb Z$ follow immediately. Inverting a
nonzero even integer kills $x$ and gives $\mathbb Q$; an odd
integer was already a unit. This proves the assertions about
all integer localizations.

The rank-two matrix
\[
 T=\begin{pmatrix}1+x&x\\x&-1+x\end{pmatrix}           \tag{C.1.19}
\]
is an involution: the off-diagonal entries of $T^2-1$ are
$2x^2$, and its diagonal entries are $2x+2x^2$ and
$-2x+2x^2$, all zero. There is a ring homomorphism
$k\to\mathbb Z/4$ sending $x$ to $2$; every element outside
$(2,x)$ maps to an odd unit, so the localization presents no
obstruction. The image of $T$ is
$\left(\begin{smallmatrix}3&2\\2&1\end{smallmatrix}\right)$.
Its Bockstein homology is all of $\mathbb F_2^2$ and its
operator is
\[
 N=\begin{pmatrix}1&1\\1&0\end{pmatrix},
 \qquad\det(tI-N)=t^2+t+1.
\]
Lemma~\ref{C:Bockstein_parity} therefore detects the nonzero
class $\eta=[(k^2,T)]-1-\mathrm{sgn}$, while it vanishes on
the degree-two polynomial image. Since $2x=0$, the two diagonal
entries of $T$ are negatives of one another. Thus
$J=\left(\begin{smallmatrix}0&1\\-1&0\end{smallmatrix}\right)$
satisfies $JT=-TJ$, and the evaluation sequence (C.1.7) gives
$2\eta=0$. The generic image is zero because $x=0$ there.

For independence in unbounded rank, take the pairwise distinct
irreducibles $f_r(t)=q_r(t^2+t)$ constructed in
Corollary~\ref{C:flat_nilpotent_infinite_kernel}, with
$d_r=\deg q_r$ and integer companion lifts $B_r$. Define
\[
 T_r=\begin{pmatrix}I_{d_r}&0\\0&-I_{d_r}\end{pmatrix}
   +x\begin{pmatrix}B_r&B_r\\B_r&B_r\end{pmatrix},
 \qquad\eta_r=[T_r]-d_r(1+\mathrm{sgn}).               \tag{C.1.20}
\]
Its square is one: the coefficient of $x$ in $T_r^2-1$ is
$\operatorname{diag}(2B_r,-2B_r)$, and the coefficient of $x^2$
has every block equal to $2B_r^2$. Under $x\mapsto2$ the
Bockstein is
\[
 N_r=\begin{pmatrix}\overline B_r&\overline B_r\\
                  \overline B_r&I+\overline B_r\end{pmatrix},
 \qquad N_r^2+N_r=\operatorname{diag}(\overline B_r,
                                     \overline B_r).
\]
Since $q_r$ is irreducible and $\overline B_r$ is a simple
$\mathbb F_2[u]$-module, the displayed identity makes
$\mathbb F_2^{2d_r}$ a module over
$\mathbb F_2[t]/(q_r(t^2+t))=\mathbb F_{2^{2d_r}}$.
Its dimension forces it to have length one over this field.
Hence $\chi_{f_s}(\eta_r)=\delta_{rs}$.
The same constant block matrix
$J=\left(\begin{smallmatrix}0&I\\-I&0\end{smallmatrix}\right)$
intertwines $T_r$ and $-T_r$, so $2\eta_r=0$ as before.
These detectors prove the asserted direct-sum injection.

At a characteristic-zero point $x=0$ and $T_r$ is the sum of
its two signs. At a characteristic-two point it has the form
$I+xU$, with $U^2=0$. Over the residue field its only simple
$C_2$-factor is trivial, so its class in the exact representation
group is its dimension, the same as $d_r(1+\mathrm{sgn})$.
Thus every $\eta_r$ is invisible at every residue-field point.
The ring is not flat over $\mathbb Z_{(2)}$, since $x\ne0$
and $2x=0$; no flatness assertion is being inferred from the
regular fibres.
\end{proof}

This also settles negatively the earlier suggested fibre-product
candidate
$\mathbb Z_{(2)}\times_{\mathbb F_2}\mathbb F_2[x]_{(x)}$:
it is the same ring, by the constant-term matching description.
In particular, regular arithmetic fibres and trivial ordinary
projective-class groups do not by themselves remove the integral
comparison kernel, even when the total ring is reduced.
